\lyseis{Β4.1}{Βασικά στερεά}
\sol{1} α) ίσα (και παράλληλα) \quad β) ύψος \quad γ) ορθογώνια \quad δ) κορυφή \quad ε) κυκλικούς \quad στ) βάση \quad ζ) σφαίρα.
\sol{2} α) κύβος \quad β) ορθογώνιο παραλληλεπίπεδο \quad γ) κύλινδρος \quad δ) κώνος \quad ε) σφαίρα \quad στ) τριγωνικό πρίσμα \quad ζ) κώνος \quad η) τετραγωνική πυραμίδα \quad θ) εξαγωνικό πρίσμα.
\sol{3} α) κύβος \quad β) κύλινδρος \quad γ) κώνος \quad δ) τετραγωνική πυραμίδα \quad ε) σφαίρα \quad στ) τριγωνικό πρίσμα.
\sol{4} α) Σωστό. \quad β) Σωστό. \quad γ) Λάθος: έχει μία βάση. \quad δ) Λάθος: είναι τρίγωνα. \quad ε) Λάθος: έχει μία βάση. \quad στ) Λάθος: δεν έχει καμία κορυφή. \quad ζ) Σωστό.
\sol{5} α) Δύο έδρες $6\times4$ (οι απέναντι), δύο έδρες $6\times3$ και δύο έδρες $4\times3$. \quad β) Τέσσερις ακμές των $6$ cm, τέσσερις των $4$ cm και τέσσερις των $3$ cm (σύνολο $12$).
\sol{6} α) Πενταγωνικό πρίσμα, $5$ παράπλευρες έδρες. \quad β) Τετραπλευρική (τετραγωνική αν είναι τετράγωνο) πυραμίδα, $4$. \quad γ) Εξαγωνικό πρίσμα, $6$. \quad δ) Τριγωνική πυραμίδα (τετράεδρο), $3$.
\sol{7} Κύβο σχηματίζουν τα α, γ και ε. Το β δεν κλείνει: τα δύο τετράγωνα που βρίσκονται πάνω από τα άκρα της σειράς πέφτουν στην ίδια έδρα. Το δ δεν κλείνει: πέντε τετράγωνα σε σειρά τυλίγονται και το πρώτο πέφτει πάνω στο πέμπτο. Το στ δεν κλείνει: διπλώνοντας, τα τετράγωνα της μιας σειράς πέφτουν πάνω στα τετράγωνα της άλλης.
\sol{8} Α: κύλινδρος. Β: τετραγωνική πυραμίδα. Γ: κώνος. Δ: τριγωνικό πρίσμα.
\sol{9} Το $\alpha$ είναι η κορυφή της πυραμίδας, το $\beta$ η βάση, το $\gamma$ το ύψος και το $\delta$ μια παράπλευρη έδρα.
\sol{10} Εξαγωνική: $7,\ 7,\ 12$. Οκταγωνική: $9,\ 9,\ 16$. Δωδεκαγωνική: $13,\ 13,\ 24$. $\nu$-γωνική: $\nu+1,\ \nu+1,\ 2\nu$. Το $\mathrm{Κ}+\mathrm{Ε}$ είναι $14,\ 18,\ 26,\ 2\nu+2$ και ισούται πάντα με $\mathrm{Α}+2$.
\sol{11} Τριγωνικό: $5,\ 6,\ 9$. Πενταγωνικό: $7,\ 10,\ 15$. Επταγωνικό: $9,\ 14,\ 21$. Δωδεκαγωνικό: $14,\ 24,\ 36$. $\nu$-γωνικό: $\nu+2,\ 2\nu,\ 3\nu$. Ισχύει $\mathrm{Κ}+\mathrm{Ε}=2\nu+\nu+2=3\nu+2=\mathrm{Α}+2$, π.χ. $6+5=11=9+2$.
\sol{12} Κουτί: $4\cdot(30+20+12)=4\cdot62=248$ cm. Κύβος: $12\cdot7=84$ cm.
\sol{13} α) $5$ έδρες: δύο ισόπλευρα τρίγωνα (οι βάσεις) και τρία ορθογώνια $4\times10$ cm. \quad β) $12\times10$ cm, γιατί η περίμετρος της βάσης είναι $3\cdot4=12$ cm. \quad γ) Ελεύθερη κατασκευή: ένα ορθογώνιο $12\times10$ χωρισμένο σε τρία ορθογώνια $4\times10$, με ένα ισόπλευρο τρίγωνο πλευράς $4$ cm κολλημένο πάνω και ένα κάτω.
\sol{14} Ελεύθερη κατασκευή. Μπροστινή όψη: ορθογώνιο $4\times6$ cm. Πάνω όψη: κύκλος ακτίνας $2$ cm. Ανάπτυγμα: ορθογώνιο $12{,}6\times6$ cm και δύο κύκλοι ακτίνας $2$ cm, κολλημένοι στις μεγάλες πλευρές του ορθογωνίου.
\sol{15} Ελεύθερη κατασκευή. Ένας κυκλικός τομέας ακτίνας $5$ cm με γωνία $144^\circ$ (μοιρογνωμόνιο) και ένας κύκλος ακτίνας $2$ cm που εφάπτεται στο τόξο του τομέα.
\sol{16} Ελεύθερη κατασκευή. Παραλληλεπίπεδο: μπροστά ορθογώνιο $4\times2$, πάνω $4\times3$, πλάι $3\times2$. Πυραμίδα: μπροστά τρίγωνο με βάση $4$ cm και ύψος $5$ cm, πάνω τετράγωνο $4\times4$ cm με τις δύο διαγωνίους του.
\sol{17} α) Κύλινδρος με ακτίνα βάσης $3$ cm και ύψος $5$ cm. \quad β) Κώνος με ακτίνα βάσης $3$ cm, ύψος $4$ cm και γενέτειρα $5$ cm. \quad γ) Σφαίρα με ακτίνα $5$ cm. Η μεγαλύτερη απόσταση δύο σημείων της είναι η διάμετρος, $10$ cm. Η μπάλα του Νίκου έχει ακτίνα $24:2=12$ cm.
\sol{18} α) $3\nu=15$, $\nu=5$: πενταγωνικό πρίσμα. \quad β) $2\nu=12$, $\nu=6$: εξαγωνική πυραμίδα. \quad γ) $2\nu=14$, $\nu=7$: επταγωνικό πρίσμα. \quad δ) Όχι, γιατί οι κορυφές του πρίσματος είναι $2\nu$, άρτιος αριθμός, και το $9$ είναι περιττός.
\sol{19} Ο Βασίλης έχει δίκιο: οι βάσεις του κύβου είναι ορθογώνια (τετράγωνα) και οι έδρες του επίσης. Η Ρήνη όχι: ένα κουτί $4\times3\times2$ είναι ορθογώνιο παραλληλεπίπεδο αλλά όχι κύβος. α) Ναι: έχει δύο παράλληλες ορθογώνιες βάσεις και ορθογώνιες παράπλευρες έδρες. \quad β) Όχι: το τριγωνικό πρίσμα δεν έχει βάσεις παραλληλόγραμμα.
\sol{20} Όχι. Η περιφέρεια της βάσης είναι $2\cdot3{,}14\cdot3=18{,}84\approx18{,}8$ cm. Το μήκος του ορθογωνίου πρέπει να είναι $18{,}8$ cm και όχι $6$ cm: η περιφέρεια είναι περίπου $3{,}14$ φορές η διάμετρος.
\sol{21} Απέναντι είναι οι $\mathrm{Α}$ και $\Gamma$, οι $\mathrm{Β}$ και $\Delta$, και οι $\mathrm{Ε}$ και $\mathrm{Ζ}$. Δίπλα στη $\mathrm{Β}$ είναι οι $\mathrm{Α}$, $\Gamma$, $\mathrm{Ε}$ και $\mathrm{Ζ}$ (όλες εκτός από τη $\Delta$ που είναι απέναντι).
\sol{22} Ελεύθερη κατασκευή. Οι όψεις του $2\times1\times1$: μπροστά ορθογώνιο $2\times1$, πάνω ορθογώνιο $2\times1$ και στο πλάι τετράγωνο $1\times1$.
\sol{23} Ελεύθερη απάντηση. Π.χ. κουτί δημητριακών (ορθογώνιο παραλληλεπίπεδο: $6$ έδρες, $12$ ακμές, $8$ κορυφές), κονσέρβα (κύλινδρος), μπάλα (σφαίρα: καμία έδρα, ακμή ή κορυφή), καπέλο πάρτι (κώνος: μία κορυφή), σκηνή (τριγωνικό πρίσμα: $5$ έδρες, $9$ ακμές, $6$ κορυφές), στέγη (πυραμίδα).
\sol{24} Ελεύθερη κατασκευή. Ο κύβος διπλώνεται σωστά αν το ανάπτυγμα είναι π.χ. σταυρός με έξι τετράγωνα $5\times5$ cm. Η πυραμίδα έχει ένα τετράγωνο $4\times4$ cm και τέσσερα ισοσκελή τρίγωνα με βάση $4$ cm και ύψος $6$ cm, ένα σε κάθε πλευρά του τετραγώνου.
